% META: {"id":"TCOMPL-CORR-CO-004","chapitre":"TCOMPL-CORRELATION-CAUSALITE","type_objet":"corrige","exercice_ref":"TCOMPL-CORR-EX-004","capacites_codes":["C5"],"sources_inspiration":[],"status":"approved","origin":"generated","status_history":["generated","audited","accepted","approved"],"release_acceptance":"RELEASE_OWNER_BATCH_ACCEPTANCE","acceptance_closure_digest":"sha256:9b3ccf9a81c5520fbb7e03b7b2d3e7bf2057a02834b6d908bf3be5f49f3b553f","acceptance_receipt_digest":"sha256:4fad86f0282e0fa4e0419cca1ad6379ae7af5a280c04a4a90880f90a1c044242"}
\begin{corrige}{TCOMPL-CORR-EX-004}
\textbf{1.} Non : ce serait absurde, les pompiers n'aggravent pas les dégâts, ils luttent contre l'incendie.

\textbf{2.} La variable cachée est la \textbf{taille de l'incendie} : un incendie plus important mobilise davantage de pompiers \textbf{et} cause davantage de dégâts. C'est la taille de l'incendie qui cause les deux variables observées, pas l'une qui cause l'autre.
\end{corrige}
